% 数学ノート集・共通テンプレート v2（積分論2026年度第1部の版面を参考）
% A5判。章内で定義・定理・例・注意等の番号を共通化。
% コンパイル： platex -interaction=nonstopmode math_note_template_v2.tex (2回)
%              dvipdfmx math_note_template_v2.dvi
\documentclass[a5paper,10pt,onecolumn]{jsarticle}
\usepackage[dvipdfmx,top=22mm,bottom=20mm,inner=19mm,outer=18mm,headheight=13pt,headsep=6mm]{geometry}
\usepackage{amsmath,amssymb,amsthm,mathtools}
\usepackage{enumitem,fancyhdr,titlesec}
\usepackage[dvipdfmx,bookmarks=true,colorlinks=false]{hyperref}
\usepackage{pxjahyper}

% ----- 差し替える文書情報 -----
\newcommand{\NoteTitle}{超実数の構成と超準解析}
\newcommand{\NoteSeries}{数学ノート集}

% ----- 文字・組版 -----
% 参考PDFの質素な明朝＋Computer Modernの印象を引き継ぎながら，
% A5判で窮屈にならないよう本文の行間を少し確保する．
\linespread{1.13}
\setlength{\parindent}{1zw}
\setlength{\parskip}{0pt}
\setlength{\abovedisplayskip}{6pt plus 2pt minus 2pt}
\setlength{\belowdisplayskip}{6pt plus 2pt minus 2pt}
\setlength{\abovedisplayshortskip}{3pt plus 2pt minus 1pt}
\setlength{\belowdisplayshortskip}{4pt plus 2pt minus 1pt}
\setlist[enumerate]{topsep=3pt,itemsep=1pt,parsep=0pt,leftmargin=2.4zw}
\titleformat{\section}{\normalfont\normalsize\bfseries}{\thesection}{0.8zw}{}
\titlespacing*{\section}{0pt}{2.0ex plus .4ex}{.9ex}
\titleformat{\subsection}{\normalfont\normalsize\bfseries}{\thesubsection}{0.7zw}{}
\titlespacing*{\subsection}{0pt}{1.5ex plus .3ex}{.6ex}
\raggedbottom

% ----- 柱とページ番号：ページ上部、罫線なし -----
\fancypagestyle{mathnote}{%
  \fancyhf{}%
  \fancyhead[L]{\footnotesize \NoteSeries\ （\NoteTitle）}%
  \fancyhead[R]{\footnotesize \thepage}%
  \renewcommand{\headrulewidth}{0pt}%
  \renewcommand{\footrulewidth}{0pt}%
}
\pagestyle{mathnote}

% ----- 節内通し番号の定義・定理群 -----
\newtheoremstyle{mathnotetheorem}
 {4.5pt}{4.5pt}{\normalfont}{}%
 {\bfseries}{}{0.65zw}%
 {\thmname{#1}\thmnumber{ #2}\thmnote{（#3）}}
\theoremstyle{mathnotetheorem}
\newtheorem{theorem}{定理}[section]
\newtheorem{definition}[theorem]{定義}
\newtheorem{lemma}[theorem]{補題}
\newtheorem{proposition}[theorem]{命題}
\newtheorem{corollary}[theorem]{系}
\newtheorem{example}[theorem]{例}
\newtheorem{remark}[theorem]{注意}
\renewcommand{\proofname}{証明}
\renewcommand{\qedsymbol}{$\square$}
\makeatletter
\renewenvironment{proof}[1][\proofname]{%
  \par\pushQED{\qed}\normalfont\topsep4\p@\@plus2\p@\relax
  \trivlist\item[\hskip\labelsep\bfseries #1]\ignorespaces
}{\popQED\endtrivlist\@endpefalse}
\makeatother

% ----- 使うノートで適宜拡張する記号 -----
\newcommand{\R}{\mathbb{R}}
\newcommand{\norm}[1]{\lVert #1\rVert}
\newcommand{\abs}[1]{\lvert #1\rvert}

\begin{document}
\section*{読了後にできること}
本稿を読むための前提は，超積，超フィルター，順序体．
本稿の具体的な到達課題は，超実数の無限小・無限大を具体的に構成して標準部分を計算する．

\section{超フィルターと超冪}
本稿では自然数集合 $\mathbb N$ 上の自由超フィルター $\mathcal U$ を固定する．
その存在には超フィルター補題を用いる．
\begin{definition}
実数列について
\[
 (x_n)\sim(y_n)\iff\{n:x_n=y_n\}\in\mathcal U
\]
と定め，商集合を
$ {}^*\mathbb R=\mathbb R^{\mathbb N}/{\sim}$ と書く．
\end{definition}
\begin{lemma}\label{lem:ultra}
任意の $A\subset\mathbb N$ について，
$A$ とその補集合のちょうど一方が $\mathcal U$ に属する．
また有限個の集合の和集合が $\mathcal U$ に属すれば，
そのうち一つが $\mathcal U$ に属する．
\end{lemma}
\begin{proof}
前半は超フィルターの極大性による．
後半が偽なら各集合の補集合が $\mathcal U$ に属する．
その共通部分は和集合の補集合に等しいので，
和集合とその補集合が共に属することとなり矛盾する．
\end{proof}
\begin{theorem}[超実数体]\label{thm:field}
$ {}^*\mathbb R$ は成分ごとの演算で体となり，
\[
 [(x_n)]<[(y_n)]\iff\{n:x_n<y_n\}\in\mathcal U
\]
により全順序体となる．
定数列は実数体を順序部分体として埋め込む．
\end{theorem}
\begin{proof}
$\sim$ の反射律・対称律は明らかであり，
推移律はフィルターの有限共通部分性から従う．
同じ理由で和・積は代表元に依存せず，
実数の環公理が商へ移る．
非零元 $x=[(x_n)]$ について
$A=\{n:x_n\ne0\}\in\mathcal U$ である．
$n\in A$ で $y_n=x_n^{-1}$，その他では $y_n=0$ とすると
$xy=1$ なので体である．
任意の二元に対する三分割
$\{x_n<y_n\},\{x_n=y_n\},\{x_n>y_n\}$ から
補題\ref{lem:ultra}を用いて三分律を得る．
順序と演算の整合性は対応する実数上の不等式を
有限個組み合わせて確認できる．
定数列が体の演算と順序を保つことは定義から従う．
\end{proof}

\section{無限小と標準部分}
\begin{definition}
$x\in{}^*\mathbb R$ が無限小であるとは
任意の実数 $\varepsilon>0$ に対して $|x|<\varepsilon$ となることをいう．
ある $M>0$ について $|x|<M$ なら有限元という．
\end{definition}
\begin{example}
$\epsilon=[(1/n)]$ は非零無限小で，
$H=[(n)]$ は任意の実数より大きい．
さらに $\epsilon H=1$ となる．
\end{example}
\begin{proof}
自由超フィルターは全ての補有限集合を含むので，
任意の $N$ について $0<\epsilon<1/N$，$H>N$ となる．
実数のArchimedes性から主張が従い，
積の等式は成分ごとの計算による．
\end{proof}
\begin{theorem}[標準部分定理]\label{thm:standard}
任意の有限元 $x$ に対し
$x-r$ が無限小となる実数 $r$ がただ一つ存在する．
これを $\operatorname{st}(x)$ と書く．
有限元 $x,y$ について
\[
 \operatorname{st}(x+y)=\operatorname{st}(x)+\operatorname{st}(y),\quad
 \operatorname{st}(xy)=\operatorname{st}(x)\operatorname{st}(y)
\]
が成り立つ．
\end{theorem}
\begin{proof}
$A=\{r\in\mathbb R:r\le x\}$ は空でなく上に有界なので
$s=\sup A$ が存在する．任意の $\varepsilon>0$ に対し
$s-\varepsilon<r\le x$ となる $r\in A$ があり，
$s+\varepsilon\notin A$ だから
$s-\varepsilon<x<s+\varepsilon$ である．
従って $x-s$ は無限小である．
二つの実数がこの条件を満たせば差の絶対値は
任意の $2\varepsilon$ より小さく，一意性が従う．
有限元の和・積は有限であり，
無限小の和および有限元と無限小の積は無限小である．
$x=r+\delta$，$y=s+\eta$ と書けば
$x+y-(r+s)=\delta+\eta$ および
$xy-rs=r\eta+s\delta+\delta\eta$ は無限小なので，
一意性から演算公式を得る．
\end{proof}

\section{L\l{}o\'sの定理と移行原理}
\begin{theorem}[L\l{}o\'sの定理]\label{thm:los}
順序体の一階言語
$\{0,1,+,\cdot,<\}$ の論理式
$\varphi(x_1,\dots,x_m)$ と
$a_j=[(a_{j,n})]$ に対し，
\[
 {}^*\mathbb R\models\varphi(a_1,\dots,a_m)
 \iff
 \{n:\mathbb R\models\varphi(a_{1,n},\dots,a_{m,n})\}
 \in\mathcal U.
\]
\end{theorem}
\begin{proof}
項についての帰納法から，超実数の項評価は
実数で成分ごとに評価した列の同値類である．
従って原子論理式の等号・不等号では定義から成立する．
論理式に関する構造帰納法を行う．
否定は超フィルターの二値性，
連言は有限共通部分性によって処理できる．
残りの命題結合子は否定と連言で表せる．
存在文が超実数で真なら証人 $b=[(b_n)]$ を取り，
帰納法の仮定から各成分における存在文の真理集合が
$\mathcal U$ に属する．
逆に実数での存在文の真理集合が $\mathcal U$ に属すれば，
その集合の各添字 $n$ で証人 $b_n$ を一つ選び，
それ以外の添字では $b_n=0$ と置く．
帰納法の仮定から $[(b_n)]$ が超実数の証人となる．
全称量化は否定と存在量化で表せるので証明が完了する．
\end{proof}
\begin{corollary}[移行原理]\label{cor:transfer}
自由変数を持たない任意の一階文 $\sigma$ について
\[
 \mathbb R\models\sigma\iff{}^*\mathbb R\models\sigma.
\]
\end{corollary}
\begin{proof}
各成分に同じ実数体を用いるので，
文の真理集合は $\mathbb N$ または空集合となる．
前者だけが $\mathcal U$ に属するから結論が従う．
\end{proof}

\section{連続性と微分可能性}
\begin{definition}
$A\subset\mathbb R$ と $f:A\to\mathbb R$ について
\[
 {}^*A=\{[(x_n)]:x_n\in A\text{ for all }n\},\quad
 {}^*f([(x_n)])=[(f(x_n))]
\]
と定める．
\end{definition}
\begin{lemma}
$ {}^*f$ は代表列に依存しない．
\end{lemma}
\begin{proof}
代表列が等しい添字集合は $\mathcal U$ に属し，
その集合上では関数値も等しいからである．
\end{proof}
\begin{theorem}[連続性の超準的特徴づけ]\label{thm:cont}
$f:A\to\mathbb R$，$a\in A$ について，
$f$ が $a$ で連続であることと，
$x\in{}^*A$，$x-a$ が無限小なら
$ {}^*f(x)-f(a)$ も無限小であることは同値である．
\end{theorem}
\begin{proof}
連続なら任意の $\varepsilon>0$ について
$\delta>0$ が存在し，
$|y-a|<\delta$ なら $|f(y)-f(a)|<\varepsilon$
となる．超実数 $x-a$ が無限小なら
$\{n:|x_n-a|<\delta\}\in\mathcal U$ なので
$|{}^*f(x)-f(a)|<\varepsilon$ となる．
逆に連続でなければ $\varepsilon_0>0$ が存在し，
$|y_n-a|<1/n$ かつ
$|f(y_n)-f(a)|\ge\varepsilon_0$ となる列を取れる．
$x=[(y_n)]$ は $a$ に無限小だけ近いが，
関数値の差は無限小ではない．
\end{proof}
\begin{theorem}[導関数の超準的特徴づけ]\label{thm:derivative}
開区間 $I$ 上の関数 $f$ と $a\in I$，
実数 $L$ について，$f'(a)=L$ であることと，
非零無限小 $h$ に対して
\[
 \frac{{}^*f(a+h)-f(a)}h-L
\]
が無限小であることは同値である．
\end{theorem}
\begin{proof}
$g(t)=(f(a+t)-f(a))/t$（$t\ne0$），
$g(0)=L$ と定める．
$f'(a)=L$ は $g$ の0での連続性と同値であり，
定理\ref{thm:cont}を適用すれば主張が従う．
\end{proof}

\section{超有限和とRiemann積分}
\begin{theorem}[超有限Riemann和]\label{thm:integral}
$f:[a,b]\to\mathbb R$ を連続関数とする．
無限大の超自然数 $H=[(N_j)]$ に対し
\[
 \operatorname{st}\left(
 \frac{b-a}{H}\sum_{k=1}^H
 {}^*f\left(a+\frac{(k-1)(b-a)}H\right)
 \right)=\int_a^bf(t)\,dt
\]
が成り立つ．
ここで超有限和は各添字で有限和を計算した列の類である．
\end{theorem}
\begin{proof}
$f$ の有界性から超有限和は有限元である．
連続関数の一様連続性によって，
通常の等分割Riemann和 $R_N$ は
$I=\int_a^bf(t)\,dt$ に収束する．
任意の $\varepsilon>0$ に対し
$N\ge N_0$ なら $|R_N-I|<\varepsilon$ とする．
$H$ が無限大なので
$\{j:N_j\ge N_0\}\in\mathcal U$ であり，
超有限和と $I$ の差は無限小である．
標準部分の一意性から表示式が従う．
\end{proof}
\begin{remark}
移行原理は一階文に対する命題である．
全ての部分集合を量化するDedekind完備性を
そのまま超実数へ移してはならない．
\end{remark}

\section{到達課題の完全解答}
\noindent\textbf{超実数の無限小の具体例．}
\begin{proposition}
非主超フィルター $\mathcal U$ による
$\mathbb R^{\mathbb N}/\mathcal U$ において
$\varepsilon=[(1/n)_{n\ge1}]$ は
正の無限小であり，
$H=[(n)_{n\ge1}]$ は正の無限大である．
しかも $\varepsilon H=1$ である．
\end{proposition}
\begin{proof}
任意の $m\ge1$ について
$\{n:0<1/n<1/m\}$ は補有限集合なので
$\mathcal U$ に属する．従って
$0<\varepsilon<1/m$ が全ての標準自然数 $m$
に対して成立する．
同様に $H>m$ となり，
各座標で $(1/n)n=1$ だから積は1である．
\end{proof}
\begin{proposition}
有限超実数 $x=2+\varepsilon$ の標準部分は2であり，
$y=1/(1+\varepsilon)$ の標準部分は1である．
\end{proposition}
\begin{proof}
$x-2=\varepsilon$ は無限小である．
また $y-1=-\varepsilon/(1+\varepsilon)$ であり
$0<\varepsilon<1$ から絶対値は $\varepsilon$ 以下である．
したがって両者とも対応する実数に無限に近い．
\end{proof}
\noindent\textbf{到達確認．}
$(1+\varepsilon)^2$ の標準部分を求めよ．
\par\noindent\textbf{解答．}
$(1+\varepsilon)^2-1=2\varepsilon+\varepsilon^2$
は無限小だから標準部分は1である．
\end{document}